\documentclass{article}
\usepackage[utf8]{inputenc}
\usepackage[german]{babel}
\usepackage{amsmath, amssymb, amsthm}
\usepackage{booktabs}
\usepackage{array}

\newcolumntype{L}[1]{>{\raggedright\arraybackslash}p{#1}}

\title{Mathematische Beschreibung endlicher Mengenfamilien}
\author{Nach Daniel J. Velleman: \emph{How to Prove It}}
\date{\today}

\begin{document}

\maketitle

\section*{Die zwei mathematischen Beschreibungen}

Wenn wir von einer ,,endlichen Familie von Teilmengen einer Topologie $\mathcal{T}$`` sprechen, können wir das auf zwei Arten formalisieren:

\subsection*{Variante A: Über eine explizite Durchnummerierung}
Wir nutzen die ersten $n$ natürlichen Zahlen als Indexmenge: $I = \{1, 2, \dots, n\}$ für ein $n \in \mathbb{N}_{\ge 1}$.
Die Familie wird geschrieben als:
\[
(U_k)_{k=1}^n \quad \text{oder} \quad \{U_1, U_2, \dots, U_n\}
\]
Die Eigenschaft, dass alle diese Mengen in der Topologie liegen, lautet in Quantorenform:
\[
\forall k \in \{1, 2, \dots, n\} \, (U_k \in \mathcal{T})
\]

\subsection*{Variante B: Über eine abstrakte endliche Indexmenge}
Wir sagen einfach, $I$ ist eine beliebige Menge, die aber die Eigenschaft hat, \emph{endlich} zu sein ($|I| < \infty$).
Die Familie wird geschrieben als:
\[
(U_i)_{i \in I}
\]
Die Eigenschaft lautet dann:
\[
\forall i \in I \, (U_i \in \mathcal{T}) \quad \land \quad |I| < \infty
\]

---

\section*{Vergleich im Scratch-Format}

Warum sind diese Schreibweisen so wichtig? Weil sie bestimmen, wie wir das \textbf{Given-Rezept für den Allquantor} anwenden dürfen. Schauen wir uns an, was passiert, wenn wir diese Familie im Beweis nutzen wollen.

\begin{center}
\small
\begin{tabular}{ L{0.45\textwidth} | L{0.45\textwidth} }
\toprule
\textbf{Nutzen von Variante A (Nummeriert)} & \textbf{Nutzen von Variante B (Abstrakt)} \\
\midrule
\textbf{Givens:} \newline
1. $\forall k \in \{1, \dots, n\} \, (U_k \in \mathcal{T})$ \newline
2. Ein konkretes $k_0 \in \{1, \dots, n\}$ existiert bereits im Beweis. &
\textbf{Givens:} \newline
1. $\forall i \in I \, (U_i \in \mathcal{T})$ \newline
2. $I$ ist endlich. \newline
3. Ein konkretes $i_0 \in I$ existiert bereits. \\
\hline
\textbf{Nachher (Instanziierung):} \newline
Du darfst die Information für dieses konkrete Element sofort nutzen: \newline
3. $U_{k_0} \in \mathcal{T}$ &
\textbf{Nachher (Instanziierung):} \newline
Du darfst die Information für dieses konkrete Element sofort nutzen: \newline
4. $U_{i_0} \in \mathcal{T}$ \\
\bottomrule
\end{tabular}
\end{center}

\vspace{1em}
\textbf{Fazit für die Praxis:}
\begin{itemize}
    \item \textbf{Variante A} ist perfekt, wenn du einen Beweis per \emph{vollständiger Induktion} führen willst (wie wir es gerade beim Durchschnitt gemacht haben), weil du dort aktiv mit der Zahl $n$ und dem Schritt zu $n+1$ arbeitest.
    \item \textbf{Variante B} ist eleganter, wenn du Eigenschaften wie ,,Stabilität unter endlichen Durchschnitten`` allgemein voraussetzen möchtest, ohne dich auf eine feste Reihenfolge festzulegen.
\end{itemize}

\end{document}
